\section{Adding oriented knots}
\label{sec:sums}
Cut a short regular arc in each of two oriented diagrams. In the sphere picture remove a small open disk about each cut, meeting each knot in a single straight interval. Glue the complementary disks, matching the outgoing end of each arc to the incoming end of the other. This is the oriented connected sum. Equivalently a box containing a knotted interval with two ends can be inserted into a short interval of the other diagram. Boundary edges may be subdivided before gluing. Both definitions give the same subdivided sphere diagram since the two end matchings are identical.

\begin{lemma}[Moving the insertion point]
An interval tangle inserted in a short regular arc can be carried to any other regular arc of the same oriented component by a finite list of planar moves. The replacement respects the orientation at the two ends.
\end{lemma}
\begin{proof}
Along a crossing-free portion use a narrow polygonal strip containing the portion and no other strands. Subdivide into boxes along the strip, and carry an oriented little box containing the tangle from one end to the other. These are planar redraws. To carry it through a crossing, straighten the carrier strand to a horizontal arc and the transverse strand to a vertical arc in the crossing disk. Put the small tangle box on the horizontal strand, with horizontal leads, wholly on the entry side of the vertical strand. Translate it to the exit side, with both leads adjusted on the horizontal line. At every intersection give the \emph{entire} carrier tangle the carrier's original height rank relative to the vertical strand.

Here is a finite movie for this translation. Choose the scan direction across the little box so that no tangle edge is parallel to a scan line, and all vertices and double points have distinct scan values (ordinary subdivisions on the leads can be omitted). A small shear, found by excluding finitely many rational slopes, achieves this, including matching the two external leads by redraws. Equivalently sweep a straight transverse line across the fixed tangle. Between two consecutive event values the intersections remain in the same order and slide along fixed segments, which is a redraw. At a regular bend traversed monotonically by the line the single crossing slides across the bend. At a local extremum the scan line creates or deletes an empty two-crossing bigon with itself over at both crossings or under at both, a type II move. At a double point of the tangle the scanning strand passes over both tangle strands or under both. The two pictures are the type III triangle with the other two ranks supplied by the original crossing. In the leads there is a single transverse crossing, which returns to its original form beyond the box. At each event choose rational times between consecutive scan values and disjoint small rational event disks; disjointness is checked by the segment tests of Section~\ref{sec:geometry}. The finitely subdivided pictures in the intervals have the same oriented plane cells. This constructs the desired list with at most one nontrivial move per bend or double point per passage. A traversal of the component passes at most $2c$ crossing visits.
\end{proof}

\begin{theorem}[The connected-sum laws]
On the oriented diagrammatic knot types, connected sum is well defined, associative and commutative and the circle is its identity. There are representatives with $c(D)+c(E)$ crossings, so this sum of crossing counts gives an upper bound for the crossing number of the sum.
\end{theorem}
\begin{proof}
The lemma moves the site in either factor. For the second factor use the symmetric gluing description with the first as the box being inserted. Exchanging the two complementary disks gives the same glued oriented sphere: one may view them as the two sides of an equator with a marked input and output. To see this in the plane, place the two cut tangles in boxes in a thin annulus, joined in their cyclic order by two seam arcs and no extra crossings. Changing which box is first is a cyclic shift, realized by rotating the annular layout through a half-turn and redrawing the boxes. If the choice of exterior face of a factor changes, place the inserted box first in the chosen small local disk of the other factor. Passing the point at infinity across a regular diagram strand can also be realized in the planar move calculus as explained below. Thus the symmetric construction gives the usual unbased connected sum. A move in either original diagram can be performed away from an insertion: slide the box to an exterior subarc first, then apply that move unchanged. A standard moving disk always has strands leaving it, so such a subarc exists. In a redraw just carry the marked interval with the redraw. Apply this construction to the given sequence of moves, with bounded induction on its prefixes.

For associativity fix three boxes with horizontal leads and join them in the indicated order. Composing the interval tangles in either parenthesization has exactly the same list of crossings and joining arcs up to a subdivision. Closing the ends gives both sides. The circle cut open is a straight unknotted interval up to redraw of its disk, so inserting it is the identity. None of these joining operations creates a new crossing in the two-box representatives.
\end{proof}

\paragraph{Changing the exterior on the sphere.}
For clarity, in the planar certificate convention a change of exterior face can be made with moves. Adjacent faces share a regular short arc. On a drawing in a large disk, replace this short arc by a detour following the boundary of the whole drawing on the side of the desired face. One explicit way is to regard the large complementary tangle as fixed in a rectangle with the two ends of the cut arc on its boundary. Take that boundary-parallel return arc from one side of the rectangle to the other by a straight line sweep across the tangle, keeping the moving arc above every strand. The moving ends are continued outside on the rectangle boundary to the same two ports. Crossing a tangle extremum or double point gives exactly the II and III events in the preceding proof. In a small disk at either port a turn past its own incident lead can also create or remove a single curl, a type I move. At the end the return arc lies in the other boundary interval, so the two sides of the cut are interchanged as exterior faces of the same embedded sphere diagram. One can obtain the rectangular cut-tangle chart by cutting the sphere around the short arc and triangulating its complementary disk with the marked ends. The redraw on each side of this construction has the same marked exterior as the corresponding endpoint. A path in the finite dual graph of the shadow changes any chosen face into any other; for a connected knot shadow it uses at most the number of faces minus one moves of this exterior-face procedure.
